% 4022 Project Thesis Template
% University of Cape Town Department of Electrical Engineering
% Template Version 0.1 - August 2021
% Adapted from the following 2 sources by J Wyngaard
% https://code.engineering.queensu.ca/ece/thesis-template-for-latex
% https://github.com/jpt13653903/LaTeX_UCT_Report.git

% -> FILE DESCRIPTION :: UCT_EE_Thesis_style.sty
%	This files defines the format of the document.

% Created in 1988 By Joseph Pallas, Stanford University
% Updated in  ??  By Bob Burke, Queen's University
% Updated in 2003 By Dean Jin, Queen's University
% Updated in 2017 By Ian Maquignaz, Queen's University
% Modified in 2021 By Jane Wyngaard, University of Cape Town





%Documentation:

%   This style file modifies the standard report or book styles to
%    follow the School of Graduate Studies and Research's "General Form
%    of Thesis."  It sets the margins and interline spacing and disallows
%   page breaks at hyphens.

%   The \beforepreface command creates the title page.  Then the user
%   should put preface section(s), using the command:

%        \prefacesection{section title}

%    The \afterpreface command then produces the tables of contents,
%    tables and figures, and sets things up to start the main body (on
%    arabic page 1).

%    The following commands can control what goes in the front matter
%    material:

%   \title{thesis title}
%    \author{author's name}
%    \supervisor{Supervisor, ie. Facilitating Professor}
%    \dept{author's department}
%        - Electrical Engineering if omitted
%    \submitdate{month year in which submitted to GPO}
%        - date LaTeX'd if omitted
%    \copyrightyear{year degree conferred}
%        - year LaTeX'd if omitted
%    \figurespagetrue or \figurespagefalse
%        - produce or don't produce a List of Figures page
%          (true by default)
%    \tablespagetrue or \tablespagefalse
%        - produce or don't produce a List of Tables page
%          (true by default)


%    The example above shows the 12pt substyle being used.  This
%    seems to give acceptable looking results, but it may be omitted
%    to get smaller print.
%}

\RequirePackage{etex}
\usepackage{fancyhdr} % For Header Formating
\usepackage{graphicx} % For Logo
\usepackage{setspace} % For \singlespacing and \doublespacing

% First thing we do is make sure that report has been loaded.  A
% common error is to try to use quthesis as a documentstyle.
\@ifundefined{chapter}{\@latexerr{The `quthesis' option should be used
with the `report' or `book' document styles}{You should probably read the
quthesis documentation.}}{}

% IK PRINTING FOR A BOOK:
% We need 1" margins except on the binding edge, where it is 1 1/2"
% Theses are single-sided, so we don't care about \evensidemargin
%\oddsidemargin 0.5in \evensidemargin 0in

\marginparwidth 40pt \marginparsep 10pt
\topmargin 0pt \headsep .5in
\textheight 8.1in \textwidth 6in

% Fixes the less than 1 inch at the bottom of the page problem
\footskip 20pt

% Disallow page breaks at hyphens (this will give some underfull vbox's,
% so an alternative is to use \brokenpenalty=100 and manually search
% for and fix such page breaks)
\brokenpenalty=10000

% %---Set up double spacing (code taken from setspace.sty)
% \ifcase \@ptsize \relax % 10pt
%     \setstretch{1.2}
%   \fi
 
% Remove the unnecessary white space around figs  

\renewcommand{\topfraction}{0.85}
\renewcommand{\textfraction}{0.1}
\renewcommand{\floatpagefraction}{0.75}

% Redefine the macro used for footnotes to use single spacing
\long\def\@footnotetext#1{\insert\footins{\setstretch{1.4}\footnotesize
    \interlinepenalty\interfootnotelinepenalty
    \splittopskip\footnotesep
    \splitmaxdepth \dp\strutbox \floatingpenalty \@MM
    \hsize\columnwidth \@parboxrestore
   \edef\@currentlabel{\csname p@footnote\endcsname\@thefnmark}\@makefntext
    {\rule{\z@}{\footnotesep}\ignorespaces
      #1\strut}}}

% Document Parameters - Set from input
\def\dept#1{\gdef\@dept{#1}}
\def\degree#1{\gdef\@degree{#1}}
\def\supervisor#1{\gdef\@supervisor{#1}}
\def\submitdate#1{\gdef\@submitdate{#1}}
\def\copyrightyear#1{\gdef\@copyrightyear{#1}}

% Document Parameters - DEFAULT PARAMETERS
\def\@title{A LaTex Thesis Template}
\def\@author{University of Cape Town}
\def\@dept{Department of Electrical Engineering}
\def\@degree{Master of Applied Science}
\def\@supervisor{A Professor}
\def\@submitdate{\ifcase\the\month\or
  January\or February\or March\or April\or May\or June\or
  July\or August\or September\or October\or November\or December\fi
  ,\space \number\the\year}
\def\@copyrightyear{\number\the\year}

% Declare If Statements for Optional Pages
\newif\ifPageFigures 
\newif\ifPageTables
\newif\ifPageGlossaries
\newif\ifPageContents
\newif\ifPageCodeSnippets
\newif\ifPageEquations

% Set Default States for Optional Pages
\PageFiguresfalse 
\PageTablesfalse
\PageGlossariesfalse
\PageContentsfalse
\PageCodeSnippetsfalse
\PageEquationsfalse

% Create the PDF
\usepackage[bookmarks,pdfauthor={\@author}, pdftitle={Thesis}]{hyperref}

% Format - Title Page
\def\titlep{
\thispagestyle{empty}
\begin{center}
    
\Huge

\vspace{1cm}

% Title
\begin{center}
\hrule 
\vskip 1mm
{\Huge \textbf{The Development of a Denoising Model for LODOX\textsuperscript{\textregistered} Statscan\textsuperscript{\textregistered} Images} }
\vskip 1mm
\hrule 
% Modify the line below to insert your subtitle.
\vskip 1mm
{\Large A Phantom Study}
\end{center}

\vspace{1mm}

% UCT logo
\includegraphics[width=0.2\linewidth]{2_Preface/uct.jpg}

\begin{center}
\large
    UNIVERSITY OF CAPE TOWN \\
    DEPT. OF ELECTRICAL ENGINEERING
\end{center}

\vspace{8mm}

% My name
\large
\textbf{Presented by:}\\
\textit{Travimadox Webb}

\vspace{8mm}
%\hline
% My supervisor's name
\textbf{Prepared for:}\\
\textit{Dr. Yaaseen Martin(Supervisor)} \\
%Department of Electrical Engineering\\
%University of Cape Town \\

\vskip 3mm
\textit{Dr. Lindie Du Plessis(Co-supervisor)}\\
%Department of Human Biology\\
%Division of Biomedical Engineering\\
%University of Cape Town


\vskip 10mm
%\hline
\begin{center}
Submitted to the Department of Electrical Engineering at the University of Cape Town in partial
fulfilment of the academic requirements for a Bachelor of Science degree in Electrical and Computer Engineering.
\end{center}

\vspace{1cm}

% Today's date
\textbf{\today}

% To add keywords?
% \vspace{2cm}
% \textbf{Keywords:} Keywords here

\end{center}

}



% Define Format - Before Preface
\def \beforepreface{
    \title
    \pagenumbering{roman}
    \pagestyle{plain}
}

% Define Format - Preface Section
\def \prefacesection#1{
    \clearpage
    \addcontentsline{toc}{chapter}{#1}
    \chapter*{#1}
}

\fancypagestyle{CustomPageStyle}{%
	\fancyhf{}% Clear header and footer
	\fancyhead[L]{\rightmark}
	\fancyhead[R]{\thepage}
	\fancyfoot[C]{\@author\space- \@dept}% Custom footer
	\renewcommand{\headrulewidth}{0.5pt}% Line at the header visible
	\renewcommand{\footrulewidth}{0.5pt}% Line at the footer visible
}

% Define Format - After Preface
\def \afterpreface{
	\clearpage
    
    % Table of Contents Page
    \ifPageContents
    	\addcontentsline{toc}{chapter}{Table of Contents}
    	\index{Table of Contents}
    	\tableofcontents
		\clearpage
	\fi
	
	% Table of Figures Page
    \ifPageFigures
        \phantomsection
        \addcontentsline{toc}{chapter}{List of Figures}
        \index{List of Figures}
        \listoffigures
        \clearpage
    \fi
    
    % Tables of Tables page
    \ifPageTables
        \phantomsection
        \addcontentsline{toc}{chapter}{List of Tables}
        \index{List of Tables}
        \listoftables
        \clearpage
    \fi
    
    % Table of Code Listings
    \ifPageCodeSnippets
        \phantomsection
	    \addcontentsline{toc}{chapter}{List of Code Listings}
	    \index{List of Code Listings}
		\lstlistoflistings
	    \clearpage
    \fi
    
    % List of Equations
    \ifPageEquations
        \phantomsection
    	\addcontentsline{toc}{chapter}{List of Equations}
    	\index{List of Equations}
	    \newfloat{Equation}{!h}{equ}[chapter]
	    \newcommand{\listequationsname}{List of Equations}
	    \listof{Equation}{\listequationsname}
	    \clearpage
    \fi
    
    % Glossary Pages
    \ifPageGlossaries
        \index{List of Symbols}
        \printglossary[type=symbol, title=List of Symbols] % Prints the acronym glossary
        \index{Nomenclature}
        \printglossary[title = Nomenclature] % Prints the terms glossary
        % \index{Glossary of Abbreviations}
        % \printglossary[type=\acronymtype,title=Glossary of Abbreviations] % Prints the abbreviation glossary
        \clearpage
    \fi
    
    \pagenumbering{arabic}    
    %\pagestyle{fancyplain}
    \pagestyle{CustomPageStyle}
}

% Redefine \thebibliography to go to a new page and put an entry in the
% table of contents
\let\@ldthebibliography\thebibliography
\renewcommand{\thebibliography}[1]{\clearpage
    \addcontentsline{toc}{chapter}{Bibliography}
    \@ldthebibliography{#1}}

